\section{Conclusion}

We present a multi-level evaluation of genotoxicity QSAR revealing that, for the Ames mutagenicity benchmark studied here, source-domain confounding arising from the mixing of drug-like and industrial chemicals drove a substantially larger performance gap ($\Delta = 0.390$ from scaffold CV to LDO) than the choice between random and scaffold splitting ($\Delta = 0.046$). Algorithm differences were small in the in-domain setting (MCC range = 0.063) but emerged under domain shift, with Random Forest showing better cross-domain robustness than gradient boosting. In~vivo micronucleus models showed ranking signal but lacked reliable classification utility with current 2D descriptors. We recommend multi-level evaluation with source-domain awareness as a minimum reporting standard for Ames mutagenicity QSAR.
